\chapter{Three-Manifolds with Positive Ricci Curvature, and Beyond}

\section{Three-manifolds with positive Ricci curvature, and beyond}

\subsection{Hamilton's theorem}

The estimates from the previous chapter show that a closed three-dimensional
manifold with positive Ricci curvature shrinks to a round point under Ricci
flow. The first singular time is finite, and after blowing up near that time one
obtains a round shrinking spherical space form.

\begin{theorem}
\label{topping-ch10-hamilton-theorem}
\dcref{10.1.1}
\group{positive-ricci-three-manifolds}
\source{topping:page-0124,topping:page-0125,topping:page-0126}
If \((\mathcal{M},g_0)\) is a closed three-dimensional Riemannian manifold
with positive Ricci curvature, then the Ricci flow \(g(t)\), with
\(g(0)=g_0\), on a maximal time interval \([0,T)\) becomes round in the
following sense. There exist
\begin{enumerate}
\item[(i)] a metric \(g_\infty\) on \(\mathcal{M}\) of constant positive sectional curvature,
\item[(ii)] a sequence \(t_i \uparrow T\),
\item[(iii)] a point \(p_\infty \in \mathcal{M}\) and a sequence \(\{p_i\} \subset \mathcal{M}\),
\end{enumerate}
such that if we define new Ricci flows \(g_i(t)\) for \(t \leq 0\) by
\[
g_i(t) = \abs{\Rm(p_i,t_i)}\, g\!\left(t_i + \frac{t}{\abs{\Rm(p_i,t_i)}}\right),
\]
then
\[
(\mathcal{M}, g_i(t), p_i) \to (\mathcal{N}, (c-t)g_\infty, p_\infty)
\]
on the time interval \((-\infty, 0]\), for some \(c > 0\).
\end{theorem}

\begin{proof}
\uses{topping-maxpr-finite-time-blowup,topping-ch8-blow-up-singularities,topping-ch9-roundness-estimate,topping-ricci-flow-uniqueness}
First, choose \(0 < \beta \leq B < \infty\) such that
\(\beta g_0 \leq \Ric(g_0) \leq B g_0\). This is possible because \(\mathcal M\)
is compact and \(\Ric(g_0) > 0\). Then \(R \geq 3\beta > 0\) at \(t = 0\), so
by Corollary~3.2.4, the maximal time \(T\) is finite.

Applying Theorem~8.5.1, we obtain a constant \(b > 0\), sequences
\(\{p_i\} \subset \mathcal M\) and \(t_i \uparrow T\), rescaled Ricci flows
\(g_i(t)\), and for \(t \in (-\infty,b)\) a limit Ricci flow
\((\mathcal N,\hat g(t))\) with base point \(p_\infty \in \mathcal N\).

By Theorem~9.7.8, for all \(t \in [0,T)\) and all \(\gamma > 0\),
\[
\left|\Ric - \frac{1}{3}Rg\right| \leq \gamma R + M(\beta,B,\gamma),
\]
with respect to \(g(t)\). Since \(g_i(t)\) is a rescaling and translation of
\(g(t)\), with respect to \(g_i(0)\) we have
\[
\left|\Ric - \frac{1}{3}Rg\right|
\leq \gamma R + \frac{M(\beta,B,\gamma)}{\abs{\Rm(p_i,t_i)}}
\]
for all \(\gamma > 0\). Passing to the limit \(i \to \infty\) gives
\[
\left|\Ric - \frac{1}{3}Rg\right| \leq \gamma R
\]
with respect to \(\hat g(0)\). Since \(\gamma > 0\) is arbitrary, it follows
that \(\Ric - \frac{1}{3}Rg = 0\) for \(\hat g(0)\). Thus \(\hat g(0)\) is
Einstein, and \(R\) is constant.

Because \(\mathcal M\) is three-dimensional, \(\hat g(0)\) has constant
sectional curvature. Since \(B > 0\) for \(g(t)\), the same holds for the
rescaled flows \(g_i(t)\), so the limit \(\hat g(0)\) has weakly positive
sectional curvature. As \(\abs{\Rm(\hat g(0))(p_\infty)} = 1\), the sectional
curvature is in fact strictly positive.

By Myers' theorem, \(\mathcal N\) is compact and hence closed. The convergence
of flows then implies \(\mathcal N = \mathcal M\). By backward uniqueness for
Ricci flow, Theorem~5.2.2, we obtain
\[
\hat g(t) = (c-t)g_\infty
\]
for some positive multiple \(g_\infty\) of \(\hat g(0)\). Therefore
\[
(\mathcal{M}, g_i(t), p_i) \to (\mathcal{M}, (c-t)g_\infty, p_\infty).
\]
\end{proof}

The existence of a metric \(g_\infty\) of constant positive sectional
curvature yields the following corollary.

\begin{corollary}
\label{topping-ch10-space-form-corollary}
\dcref{10.1.2}
\group{positive-ricci-three-manifolds}
\source{topping:page-0126}
\uses{topping-ch10-hamilton-theorem}
Any closed Riemannian three-manifold with positive Ricci curvature admits a
metric of constant positive sectional curvature. In particular, if the manifold
is also simply connected, then it must be \(S^3\).
\end{corollary}

\subsection{Beyond the case of positive Ricci curvature}

To move beyond the assumption \(\Ric > 0\), one studies singularities and
surgeries. The first key input is a curvature pinching theorem of Hamilton and
Ivey.

\begin{theorem}
\label{topping-ch10-hamilton-ivey-theorem}
\dcref{10.2.1}
\group{positive-ricci-three-manifolds}
\source{topping:page-0126}
\uses{topping-ch9-ode-pde-theorem}
For all \(B < \infty\) and \(\varepsilon > 0\), there exists
\(M = M(\varepsilon,B) < \infty\) such that if \((\mathcal{M}, g_0)\) is a
closed three-manifold with \(\abs{\Rm(g_0)} \leq B\), then the Ricci flow
\(g(t)\) with \(g(0) = g_0\) satisfies
\[
E + (\varepsilon R + M)g \geq 0,
\]
where \(E\) is the Einstein tensor.
\end{theorem}

This means that the lowest sectional curvature satisfies
\[
\lambda_1 \geq -\varepsilon R - M,
\]
and since \(R = 2(\lambda_1 + \lambda_2 + \lambda_3)\), a very negative
sectional curvature forces the most positive sectional curvature to be much
larger still.

\begin{corollary}
\label{topping-ch10-blowup-weakly-positive}
\dcref{10.2.2}
\group{positive-ricci-three-manifolds}
\source{topping:page-0127}
\uses{topping-ch10-hamilton-ivey-theorem}
For a Ricci flow on a closed three-dimensional manifold, when we take a
blow-up limit
\[
(\mathcal{M}, g_i(t), p_i) \to (\mathcal{N}, \hat{g}(t), p_\infty),
\]
as in Theorem~8.5.1, the manifold \((\mathcal{N}, \hat{g}(t))\) has weakly
positive sectional curvature for all \(t \in (-\infty, b)\).
\end{corollary}

\begin{remark}
\label{topping-ch10-perelman-reference}
\dcref{10.2.3}
\group{positive-ricci-three-manifolds}
\source{topping:page-0127}
This severely constrains the singularities which are possible. To understand
more about the Ricci flow beyond these notes, the reader is directed first to
Perelman's paper.
\end{remark}
